% SPDX-License-Identifier: CC-BY-SA-4.0
% Author: Matthieu Perrin

\begingroup

\begin{exercice}[Description d'expressions rationnelles] \label{exo:lexing/regexp/read_regexp}

  Décrire les langages définis par les expressions rationnelles suivantes.\\
  La notation $[c_1\!-\!c_2]$ désigne un symbole parmi tous les symboles entre $c_1$ et $c_2$.
  \begin{question}
  \item $([1\!-\!9][0\!-\!9]^\star (0\mid 2\mid 4\mid 6\mid 8)) \mid  (0\mid 2\mid 4\mid 6\mid 8)$     \hspace\fill sur l'alphabet $\Sigma_1 = \{0,1,2,3,4,5,6,7,8,9,H\}$.
  \item $([0\!-\!9]\mid 10\mid 11)H[0\!-\!5][0-9]$                        \hspace\fill sur l'alphabet $\Sigma_1 = \{0,1,2,3,4,5,6,7,8,9,H\}$.
  \item $(a\mid b\mid c)^\star  a(a\mid b\mid c)^\star  a(a\mid b\mid c)^\star $             \hspace\fill sur l'alphabet $\Sigma_3 = \{a,b,c\}$.
  \end{question}

\end{exercice}

\begin{correction}

  \begin{question}
  \item L'écriture décimale des nombres pairs, sans zéro intuile.
  \item Les heures sur 12 heures de la forme hhHmm.
  \item Tous les mots $u$ sur $\Sigma_3$ contenant au moins deux `a' pas forcément consécutifs.
  \end{question}

\end{correction}

\endgroup
\endinput
